\documentclass{article}

\usepackage{amsmath, amsthm, amssymb, amsfonts}
\usepackage{thmtools}
\usepackage{graphicx}
\usepackage{setspace}
\usepackage{geometry}
\usepackage{float}
\usepackage[hidelinks]{hyperref}
\usepackage[utf8]{inputenc}
\usepackage[spanish]{babel}
\usepackage{framed}
\usepackage[dvipsnames]{xcolor}
\usepackage{tcolorbox}
\usepackage{tikz}
\usepackage{caption}
\usepackage{longtable}
\usepackage{pdflscape}
\usepackage{svg}
\usepackage{subcaption}
\usepackage{caption}
\usepackage{multirow}
\usepackage{array}
\usepackage{listings}
\usepackage{cancel}
\usepackage{fancyhdr}


\colorlet{LightGray}{White!90!Periwinkle}
\colorlet{LightOrange}{Orange!15}
\colorlet{LightGreen}{Green!15}



\newcommand{\HRule}[1]{\rule{\linewidth}{#1}}

\declaretheoremstyle[name=Theorem,]{thmsty}
\declaretheorem[style=thmsty,numberwithin=section]{theorem}
\tcolorboxenvironment{theorem}{colback=LightGray}

\declaretheoremstyle[name=Proposition,]{prosty}
\declaretheorem[style=prosty,numberlike=theorem]{proposition}
\tcolorboxenvironment{proposition}{colback=LightOrange}

\declaretheoremstyle[name=Principle,]{prcpsty}
\declaretheorem[style=prcpsty,numberlike=theorem]{principle}
\tcolorboxenvironment{principle}{colback=LightGreen}

\newcolumntype{L}[1]{>{\raggedleft\let\newline\\\arraybackslash\hspace{0pt}}m{#1}}
\newcolumntype{C}[1]{>{\centering\let\newline\\\arraybackslash\hspace{0pt}}m{#1}}
\newcolumntype{R}[1]{>{\raggedright\let\newline\\\arraybackslash\hspace{0pt}}m{#1}}

\setstretch{1.2}
\geometry{
    textheight=9in,
    textwidth=5.5in,
    top=1in,
    headheight=12pt,
    headsep=25pt,
    footskip=30pt
}

\lstdefinestyle{bashstyle}{
    language=bash,
    basicstyle=\ttfamily,
    backgroundcolor=\color{gray!10},
    keywordstyle=\color{blue},
    commentstyle=\color{green!40!black},
    stringstyle=\color{red},
    showstringspaces=false,
    numbers=left,
    numberstyle=\tiny\color{gray},
    breaklines=true,
    breakatwhitespace=true,
    frame=tb,
    rulecolor=\color{black!70},
    framerule=0.5pt,
    tabsize=4,
    captionpos=b
}

\lstdefinestyle{javastyle}{
    language=Java,
    basicstyle=\ttfamily,
    backgroundcolor=\color{gray!10},
    keywordstyle=\color{blue},
    commentstyle=\color{green!40!black},
    stringstyle=\color{red},
    showstringspaces=false,
    numbers=left,
    numberstyle=\tiny\color{gray},
    breaklines=true,
    breakatwhitespace=true,
    frame=tb,
    rulecolor=\color{black!70},
    framerule=0.5pt,
    tabsize=4,
    captionpos=b
}

\lstdefinestyle{pythonstyle}{
    language=Python,
    basicstyle=\ttfamily,
    backgroundcolor=\color{gray!10},
    keywordstyle=\color{blue},
    commentstyle=\color{green!40!black},
    stringstyle=\color{red},
    showstringspaces=false,
    numbers=left,
    numberstyle=\tiny\color{gray},
    breaklines=true,
    breakatwhitespace=true,
    frame=tb,
    rulecolor=\color{black!70},
    framerule=0.5pt,
    tabsize=4,
    captionpos=b
}

\lstdefinestyle{cppstyle}{
    language=C++,
    basicstyle=\ttfamily,
    backgroundcolor=\color{gray!10},
    keywordstyle=\color{blue},
    commentstyle=\color{green!40!black},
    stringstyle=\color{red},
    showstringspaces=false,
    numbers=left,
    numberstyle=\tiny\color{gray},
    breaklines=true,
    breakatwhitespace=true,
    frame=tb,
    rulecolor=\color{black!70},
    framerule=0.5pt,
    tabsize=4,
    captionpos=b
}

% --------------------------------------------------------------------
% TITLE PAGE
% --------------------------------------------------------------------
\title{
    \vspace{2cm}
    \HRule{1pt} \\[0.5cm]
    \textbf{\Huge Plantilla de Programación Competitiva}\\[0.3cm]
    \HRule{1pt} \\[1.5cm]
    \LARGE Universidad de Bogotá Jorge Tadeo Lozano \\[0.3cm]
    \Large Facultad de Ciencias Naturales e Ingeniería\\[0.3cm]
    \Large Semillero de Programación Competitiva\\[1.5cm]
    \large \textbf{Maratón Nacional ACIS REDIS 2025}\\[2cm]
}

\author{
    \Large Elaborado por: Ludwig Alvarado y Nicolás Leitón
}

\date{\vspace{1cm} \Large Bogotá D.C., \today}


% --------------------------------------------------------------------
\begin{document}

\maketitle
\newpage

\tableofcontents

\thispagestyle{empty}
\newpage
\setcounter{page}{1}




\section{Leer Input}

\subsection{Python}

\subsubsection{n inicial y luego casos}

\begin{lstlisting}[style=pythonstyle]
import sys
inputs = iter(sys.stdin.readlines())
TC = int (next(inputs))
for _ in range(TC):
  print(sum(map(int, next(inputs).split())))

\end{lstlisting}

\begin{lstlisting}
3
1 2
5 7
6 3
--- out
3
12
9
\end{lstlisting}

\subsubsection{Terminar cuando es 0 0}

\begin{lstlisting}[style=pythonstyle]
import sys
for line in sys.stdin.readlines():
  if line=='0 0\n': break
  print(sum(map(int, line.split())))
\end{lstlisting}

\begin{lstlisting}
1 2
5 7
6 3
0 0
1 1 <- no se procesa
--- out
3
12
9
\end{lstlisting}


\subsubsection{Enumerar casos hasta EOF}

\begin{lstlisting}[style=pythonstyle]

import  sys
for c, line in enumerate(sys.stdin.readlines(), 1):
    print("Case %s: %s\n" % (c, sum(map(int, line.split()))))

\end{lstlisting}


\begin{lstlisting}
1 2
5 7
6 3
--- out
Case 1: 3

Case 2: 12

Case 3: 9

\end{lstlisting}

\subsubsection{Input matriz nxm}

\begin{lstlisting}[style=pythonstyle]

import sys
lines = sys.stdin.readlines()
i= 0
matrices = []

while i<len(lines):
  # Leer dimensiones
  n, m = map(int, lines[i].split())
  i += 1

  #Leer la matriz de tamano n
  matrix = [list(map(int, lines[i+j].split())) for j in range(n)]
  matrices.append(matrix)

  i +=n

\end{lstlisting}




\subsection{C++}

\subsubsection{n inicial y luego casos}

\begin{lstlisting}[style=cppstyle]

int n; cin >> n;
while(n--){
  int a, b; cin >> a >> b;
  cout << (a + b) << endl;
}

\end{lstlisting}



\begin{lstlisting}
3
1 2
5 7
6 3
--- out
3
12
9
\end{lstlisting}


\subsubsection{Terminar cuando es 0 0}

\begin{lstlisting}[style=cppstyle]
int a, b;
while(cin >> a >> b, (a || b))
  cout << (a + b) << endl;

\end{lstlisting}

\begin{lstlisting}
1 2
5 7
6 3
0 0
1 1 <- no se procesa
--- out
3
12
9
\end{lstlisting}

\subsubsection{Leer líneas hasta EOF}

\begin{lstlisting}[style=cppstyle]
string line;
while(getline(cin, line)){
  stringstream ss(line);
  int x, sum = 0;
  while(ss >> x) sum += x;
  cout << sum << endl;
}
\end{lstlisting}

\begin{lstlisting}
3 4 5
10 20
7 8 9 10
--- out
12
30
34

\end{lstlisting}

\subsubsection{Leer N líneas (string)}

\begin{lstlisting}[style=cppstyle]
int N;
if (!(cin >> N)) return 0;
cin.ignore();
for(int i = 0; i < N; ++i){
  string s;
  getline(cin, s);
}
\end{lstlisting}


\subsubsection{Leer vector}

\begin{lstlisting}[style=pythonstyle]
int n; cin >> n;

vector<int> v(n);
copy_n(istream_iterator<int>(cin), n, v.begin());

\end{lstlisting}





\section{Output Format}


\subsection{Python}

\begin{lstlisting}[style=pythonstyle]
#Decimals output
text = "The price is {:.2f} dollars"
print(txt.format(45))
# The price is 45.00 dollar

#Use ',' as thousand separator:
txt = "the universe is {:,} years old "
print(txt.format(13800000000))
# The universe is 13,800,000,000 years old

# Convert to percentage format:
txt ="you scored {:.1%}"
print(txt.format(0.255)''
# you scored 25.5%

txt ="you scored {:.0%}"
print(txt.format(0.255)''
# you scored 26%


#covert number to binary

txt = "The binary version of {0} is {0:b}"
print(txt.format(5))
#The binary version of 5 is 101



\end{lstlisting}

\subsection{C++}

\begin{lstlisting}[style=cppstyle]

vector<int> v = {1, 2, 3, 4, 5};
copy(v.begin(), v.end(), ostream_iterator<int>(cout));
// 12345%

double x = 123.456789;
cout << fixed << setprecision(3);
cout << "Value: " << x << endl;
// Value: 123.457

cout << setw(10) << setfill('*') << 42 << endl;
// **********42

cout << boolalpha << (x > 100) << endl;
cout << hex << 255 << endl;
// true
// ff

vector<vector<int>> mat = {{1, 2, 3}, {4, 5, 6}, {7, 8, 9}};
for(auto &row : mat){
  for (int i = 0; i < row.size(); ++i)
    cout << (i ? " " : "") << row[i];
  cout << endl;
}
/*
1 2 3
4 5 6
7 8 9
*/
\end{lstlisting}


\section{Macros}

\subsection{Python}

\begin{lstlisting}[style=pythonstyle]
def crear_matriz(n, m, valor =0):
  return [[valor for _ in range(m)] for _ in range(n)]

def imprimir_matriz(mat):
  for fila in mat:
     print(*fila)

n, m = 3, 4
matriz= crear_matriz(n, m)
imprimir_matriz(matriz)

"""
0 0 0 0
0 0 0 0
0 0 0 0
"""

\end{lstlisting}

\subsection{C++}


\begin{lstlisting}[style=cppstyle]
#include <bits/stdc++.h>
using namespace std;
#define fast ios::sync_with_stdio(false); cin.tie(nullptr); cout.tie(nullptr)
#define endl '\n'
#define pb push_back
#define all(x) (x).begin(), (x).end()
#define sz(x) (int)(x).size()
#define for0(i, n) for (int i = 0; i < (int)(n); ++i) // base 0 index
#define for1(i, n) for (int i = 0; i < (int)(n); ++i) // base 1 index
#define forc(i, l, r) for (int i = (int)(l); i <= (int)(r); ++i) // intervalo cerrado de l a r inclusivo
#define forr0(i, n) for (int i = (int)(n) - 1; i >= 0 ; --i) // al contrario hasta 0.
#define forr1(i, n) for (int i = (int)(n) - 1; i > 0; --i) // al contrario hasta 1

typedef long long ll;
typedef short int si;
typedef unsigned long long ull;
typedef long double ld;
typedef unsigned int ui;
typedef string str;
typedef pair<int, int> pii;
typedef pair<ll, ll> pll;
typedef vector<int> vi;


const int INFI = INT_MAX;
const long INFL = LONG_MAX;
const long long INFLL = LLONG_MAX;

\end{lstlisting}


\section{Operaciones Bitwise}

\subsection{C++}

\begin{lstlisting}[style=cppstyle]

// Macros para operaciones bitwise
#define isOn(S, j) (S & (1 << j))
#define setBit(S, j) (S |= (1 << j))
#define clearBit(S, j) (S &= ~(1 << j))
#define toggleBit(S, j) (S ^= (1 << j))
#define lowBit(S) (S & (-S))
#define setAll(S, n) (S = (1 << n)-1)

#define modulo(S, N) ((S) & (N-1)) // Retorna S % N, donde N es una potencia de 2
#define isPowerOfTwo(S) (!(S) & (S-1))
#define nearestPowerOfTwo(S) (1 << lround(log2(S)))
#define turnOffLastBit(S) ((S) & (S-1))
#define turnOnLastZero(S) ((S) | (S+1))
#define turnOffLastConsecutiveBits(S) ((S) & (S+1))
#define turnOnLastConsecutiveZeroes(S) ((S) | (S-1))

\end{lstlisting}


\section{Bitmask y Binarios}

\begin{lstlisting}[style=cppstyle]
  ll binpow(ll a, ll b){
    ll res = 1;
    ll counter = 0;
    while (b > 0){
      bitset<16> binary(b);
      if ( b& 1 ) res = res * a;
      a + a * a;
      counter++;
      b >>=1;
    }
    return res;
  }

  ll binpowmod(ll a, ll b, ll m){
    a \%= m;
    ll res = 1;
    while ( b > 0 ){
      if ( b & 1 ) res = res * a \% m;
      a = a * a \% m;
      b >>= 1;
    }
    return res;
  }


  bitset<8> b1; // 00000000
  bitset<8> b2(5); // 00000101
  bitset<8> b3(string("1101")); // 00001101
  bitset<8> b("10110010")
  b[0] // 0 (least significant bit)
  b[7] // 1 (most significant bit)
  b[0] = 1; // Set bit 0 10110011
  b.set(3); // set bit at index 3 10111011
  b.reset(1); // reset bit at index 1
  b.flip(2); // toggle bit at index 2 10111101
  b.flip();

  b.any() // True if any bit == 1
  b.none() // true if all bits == 0
  b.all() // true if all bits == 1
  b.count() // number of bits == 1
  b.size() // number of bits in bitset

  b.to_ulon() // unsigned long
  b.to_ullong() // unsigned long long
  b.to_string() // string

// Obtener los subconjuntos de una mascara en binario
vector<int> subset_bitmask(int mask){
    vector<int> subsets;
    for (int subset = mask; subset; subset = (mask & (subset -1))) subsets.push_back(subset);
    return subsets;
}
\end{lstlisting}

\section{Swap dos números}


\begin{lstlisting}[style=cppstyle]
void swap_two_iters(vector<int> vec, int n){
    const auto itr = find(all(vec), n);
    auto indx = distance(vec.begin(), itr);
    iter_swap(vec.begin() + n, itr);
    iter_swap(vec.begin() + n, vec.begin() + indx);
}

int num = 2;
const auto itr = find(all(v), num);
auto indx = distance(v.begin(), itr);
iter_swap(v.begin() + num, v.begin() + indx);
\end{lstlisting}


\section{Grafos: DFS, BFS y Dijkstra}

Esta sección contiene implementaciones base y trucos para grafos no dirigidos o dirigidos, con pesos o sin ellos.

\subsection{Representación del Grafo}

\begin{lstlisting}[style=cppstyle]
// Grafo con N nodos numerados 0..N-1
int N;
vector<vector<int>> adj(N); // sin peso
vector<vector<pair<int,int>>> g(N); // con pesos (v, w)

for(int i=0;i<M;i++){
    int u,v; cin>>u>>v;
    adj[u].push_back(v);
    adj[v].push_back(u); // quitar si es dirigido
}
\end{lstlisting}


\subsection{DFS (Depth First Search)}

Complejidad: \(O(N + M)\)

\begin{lstlisting}[style=cppstyle]
vector<vector<int>> adj;
vector<int> vis;

void dfs(int u){
    vis[u] = 1;
    for(int v : adj[u])
        if(!vis[v])
            dfs(v);
}

int main(){
    int n,m; cin>>n>>m;
    adj.assign(n,{});
    vis.assign(n,0);
    while(m--){
        int u,v; cin>>u>>v;
        adj[u].push_back(v);
        adj[v].push_back(u);
    }
    dfs(0); // desde nodo 0
}
\end{lstlisting}

\textbf{Trucos:}
\begin{itemize}
  \item Útil para contar componentes conexas.
  \item Para detectar ciclos: pasar padre y verificar si hay arista a nodo visitado distinto del padre.
  \item Usar iterativo con \lstinline|stack| si hay riesgo de stack overflow.
\end{itemize}


\subsection{BFS (Breadth First Search)}

Complejidad: \(O(N + M)\)

\begin{lstlisting}[style=cppstyle]
vector<vector<int>> adj;
vector<int> dist;

void bfs(int s){
    queue<int> q;
    dist.assign(adj.size(), -1);
    dist[s] = 0;
    q.push(s);

    while(!q.empty()){
        int u = q.front(); q.pop();
        for(int v : adj[u])
            if(dist[v] == -1){
                dist[v] = dist[u] + 1;
                q.push(v);
            }
    }
}
\end{lstlisting}

\textbf{Trucos:}
\begin{itemize}
  \item Útil para hallar distancia mínima en grafos no ponderados.
  \item Para reconstruir camino: mantener vector \lstinline|parent[v]|.
\end{itemize}


\subsection{Dijkstra (Shortest Path con pesos positivos)}

Complejidad: \(O((N + M) \log N)\) con \lstinline|priority_queue|.

\begin{lstlisting}[style=cppstyle]
vector<vector<pair<int,int>>> g; // (v, w)
const ll INF = 1e18;

vector<ll> dijkstra(int s){
    vector<ll> dist(g.size(), INF);
    priority_queue<pair<ll,int>, vector<pair<ll,int>>, greater<pair<ll,int>>> pq;
    dist[s] = 0;
    pq.push({0,s});

    while(!pq.empty()){
        auto [du, u] = pq.top(); pq.pop();
        if(du != dist[u]) continue;
        for(auto [v, w] : g[u]){
            if(dist[v] > du + w){
                dist[v] = du + w;
                pq.push({dist[v], v});
            }
        }
    }
    return dist;
}
\end{lstlisting}

\textbf{Trucos:}
\begin{itemize}
  \item Si hay muchas aristas (\(M \approx N^2\)), usar vector de mínimos (\(O(N^2)\)).
  \item Usar \lstinline|greater<pair<ll,int>> pq;| para min-heap.
  \item Para camino mínimo: mantener vector \lstinline|parent[v]|.
\end{itemize}


\subsection{Extras y Trucos de Grafos}

\textbf{1. Componentes Conexas:}
\begin{lstlisting}[style=cppstyle]
int comp = 0;
for(int i=0;i<n;i++)
  if(!vis[i]){ dfs(i); comp++; }
\end{lstlisting}

\textbf{2. BFS 0–1 (para grafos con pesos 0 y 1):}
\begin{lstlisting}[style=cppstyle]
deque<int> dq;
vector<int> dist(n, 1e9);
dist[s] = 0;
dq.push_back(s);

while(!dq.empty()){
    int u = dq.front(); dq.pop_front();
    for(auto [v,w]: g[u]){
        if(dist[v] > dist[u] + w){
            dist[v] = dist[u] + w;
            if(w==1) dq.push_back(v);
            else dq.push_front(v);
        }
    }
}
\end{lstlisting}

\textbf{3. Detectar Ciclos en Dirigido (Toposort Check):}
\begin{lstlisting}[style=cppstyle]
vector<int> indeg(n);
for(int u=0;u<n;u++)
  for(int v:adj[u]) indeg[v]++;
queue<int> q;
for(int i=0;i<n;i++) if(!indeg[i]) q.push(i);
int cnt=0;
while(!q.empty()){
  int u=q.front(); q.pop(); cnt++;
  for(int v:adj[u]) if(--indeg[v]==0) q.push(v);
}
bool has_cycle = (cnt < n);
\end{lstlisting}

\textbf{4. BFS multi-source:}
\begin{lstlisting}[style=cppstyle]
queue<int> q;
for(int s: fuentes){
    dist[s]=0;
    q.push(s);
}
\end{lstlisting}

\textbf{5. Bidirectional BFS (más rápido si inicio y fin conocidos).}

\textbf{6. Usar \lstinline|vector<vector<pair<int,int>>>| si los pesos son necesarios.}

\textbf{7. Reservar memoria antes:} \lstinline|adj.assign(n, {});|

---

\textbf{Resumen de Complejidades:}

\begin{center}
\begin{tabular}{lcc}
\hline
Algoritmo & Complejidad & Uso principal \\
\hline
DFS & \(O(N+M)\) & Conectividad, topología, ciclos \\
BFS & \(O(N+M)\) & Distancias en grafos no ponderados \\
Dijkstra & \(O((N+M)\log N)\) & Caminos mínimos (pesos positivos) \\
0–1 BFS & \(O(N+M)\) & Caminos mínimos (pesos 0/1) \\
TopoSort (Kahn) & \(O(N+M)\) & DAG orden topológico \\
\hline
\end{tabular}
\end{center}




\section{Estructuras de Datos Built-In}

Esta sección contiene las estructuras y algoritmos más útiles de la STL para programación competitiva (C++11).  
Incluye las operaciones más comunes y su complejidad temporal.

\subsection{Vector (\texttt{std::vector})}

Array dinámico.

\begin{itemize}
  \item Acceso por índice: \(O(1)\)
  \item Inserción/Eliminación al final: \(O(1)\) amortizado
  \item Inserción/Eliminación en medio: \(O(n)\)
  \item Redimensionar: \(O(n)\)
\end{itemize}

\begin{lstlisting}[style=cppstyle]
vector<int> v;
v.push_back(10);
v.pop_back();
v.resize(5);
v.clear();
v.reserve(1000); // evita realocaciones
int x = v.back();
\end{lstlisting}


\subsection{Set y Multiset (\texttt{std::set}, \texttt{std::multiset})}

Árbol balanceado (Red-Black Tree).

\begin{itemize}
  \item Insertar/Eliminar/Buscar: \(O(\log n)\)
  \item lower\_bound / upper\_bound: \(O(\log n)\)
\end{itemize}

\begin{lstlisting}[style=cppstyle]
set<int> s;
s.insert(3);
s.insert(1);
s.insert(2);
s.erase(1);
if(s.find(3) != s.end()) cout << "found\n";

for(int x : s) cout << x << " "; // 2 3
\end{lstlisting}

\textbf{multiset} permite duplicados y mantiene orden.

\begin{lstlisting}[style=cppstyle]
multiset<int> ms = {1, 1, 2, 3};
ms.erase(ms.find(1)); // borra solo una ocurrencia
\end{lstlisting}


\subsection{Map y Multimap (\texttt{std::map}, \texttt{std::multimap})}

Diccionario ordenado por clave.

\begin{itemize}
  \item Insertar/Eliminar/Buscar: \(O(\log n)\)
  \item Acceso por clave: \(O(\log n)\)
\end{itemize}

\begin{lstlisting}[style=cppstyle]
map<string, int> freq;
freq["apple"]++;
freq["banana"] = 3;

for (auto &p : freq)
  cout << p.first << ": " << p.second << endl;
\end{lstlisting}


\subsection{Unordered Map / Set (\texttt{std::unordered\_map}, \texttt{std::unordered\_set})}

Tablas hash (no ordenadas).

\begin{itemize}
  \item Inserción/Eliminación/Búsqueda promedio: \(O(1)\)
  \item Peor caso: \(O(n)\) (colisiones)
\end{itemize}

\begin{lstlisting}[style=cppstyle]
unordered_map<string,int> mp;
mp["dog"] = 5;
if(mp.count("dog")) cout << mp["dog"];
\end{lstlisting}

\textbf{Recomendado} para conteos rápidos o lookups cuando el orden no importa.


\subsection{Deque (\texttt{std::deque})}

Cola doble: acceso eficiente en ambos extremos.

\begin{itemize}
  \item push\_front / push\_back: \(O(1)\)
  \item pop\_front / pop\_back: \(O(1)\)
  \item Acceso por índice: \(O(1)\)
\end{itemize}

\begin{lstlisting}[style=cppstyle]
deque<int> dq;
dq.push_back(1);
dq.push_front(2);
cout << dq.front() << " " << dq.back(); // 2 1
\end{lstlisting}


\subsection{Stack / Queue / Priority Queue}

\textbf{Stack:} LIFO (\(O(1)\) push/pop/top)

\begin{lstlisting}[style=cppstyle]
stack<int> st;
st.push(5);
cout << st.top(); // 5
st.pop();
\end{lstlisting}

\textbf{Queue:} FIFO (\(O(1)\) push/pop/front)

\begin{lstlisting}[style=cppstyle]
queue<int> q;
q.push(1);
q.push(2);
q.pop();
cout << q.front(); // 2
\end{lstlisting}

\textbf{Priority Queue:} heap binario (\(O(\log n)\))

\begin{lstlisting}[style=cppstyle]
priority_queue<int> pq; // max-heap
pq.push(10);
pq.push(5);
cout << pq.top(); // 10
pq.pop();

priority_queue<int, vector<int>, greater<int>> minpq; // min-heap
\end{lstlisting}


\subsection{Pair y Tuple}

Emparejar valores o agrupar varios tipos.

\begin{lstlisting}[style=cppstyle]
pair<int, string> p = {1, "abc"};
auto [x, y] = p;

tuple<int,int,string> t = {1, 2, "hi"};
auto [a, b, c] = t;
\end{lstlisting}


\subsection{Algoritmos Comunes (\texttt{<algorithm>})}

\begin{itemize}
  \item \lstinline|sort(all(v))| → \(O(n \log n)\)
  \item \lstinline|reverse(all(v))| → \(O(n)\)
  \item \lstinline|binary_search(all(v), x)| → \(O(\log n)\)
  \item \lstinline|lower_bound(all(v), x)| → \(O(\log n)\)
  \item \lstinline|max_element(all(v))| → \(O(n)\)
  \item \lstinline|accumulate(all(v), 0)| → \(O(n)\)
\end{itemize}

\begin{lstlisting}[style=cppstyle]
vector<int> v = {3, 1, 4, 2};
sort(all(v));
reverse(all(v));
if(binary_search(all(v), 3)) cout << "found";
\end{lstlisting}


\subsection{Bitset (\texttt{std::bitset<N>})}

Arreglo de bits estático.

\begin{itemize}
  \item Acceso/Set/Reset: \(O(1)\)
  \item Operaciones bitwise: \(O(N / \text{word\_size})\)
\end{itemize}

\begin{lstlisting}[style=cppstyle]
bitset<8> b("10100110");
b.set(3);
b.reset(1);
b.flip();
cout << b.count(); // num de bits encendidos
\end{lstlisting}


\subsection{String (\texttt{std::string})}

\begin{itemize}
  \item Acceso: \(O(1)\)
  \item Comparar: \(O(\min(n,m))\)
  \item Substring: \(O(k)\)
\end{itemize}

\begin{lstlisting}[style=cppstyle]
string s = "hello world";
cout << s.substr(0,5);
reverse(all(s));
if(s.find("he") != string::npos) cout << "found";
\end{lstlisting}


\subsection{Union-Find / Disjoint Set (DSU)}

Estructura esencial para componentes conectadas o MST.  
Complejidad amortizada \(O(\alpha(n)) \approx O(1)\).

\begin{lstlisting}[style=cppstyle]
struct DSU {
    vector<int> p, sz;
    DSU(int n): p(n), sz(n,1){ iota(p.begin(), p.end(), 0); }
    int find(int x){ return p[x]==x ? x : p[x]=find(p[x]); }
    bool unite(int a, int b){
        a=find(a); b=find(b);
        if(a==b) return false;
        if(sz[a]<sz[b]) swap(a,b);
        p[b]=a; sz[a]+=sz[b];
        return true;
    }
};
\end{lstlisting}


\subsection{Recomendaciones Finales}

\begin{lstlisting}[style=cppstyle]
// Siempre incluir
ios::sync_with_stdio(false);
cin.tie(nullptr);

// Alias y macros utiles
using ll = long long;
using pii = pair<int,int>;
#define all(x) (x).begin(), (x).end()
#define rall(x) (x).rbegin(), (x).rend()
#define sz(x) int((x).size())

// RNG
mt19937 rng(chrono::steady_clock::now().time_since_epoch().count());
\end{lstlisting}

\section{Expresiones Regulares (Regex) en C++}

La biblioteca \lstinline|<regex>| permite realizar coincidencias de patrones en cadenas usando expresiones regulares, útiles para validación, búsqueda y extracción de información.

\subsection{Incluir y estructuras básicas}
\begin{lstlisting}[style=cppstyle]
#include <bits/stdc++.h>
using namespace std;

int main() {
    string s = "abc123xyz";
    regex re("[0-9]+"); // uno o mas digitos
    if (regex_search(s, re))
        cout << "Tiene numeros!\n";
}
\end{lstlisting}

\subsection{Principales clases y funciones}

\begin{itemize}
  \item \lstinline|regex re("pattern");| — Compila un patrón.
  \item \lstinline|regex_match(str, re)| — Comprueba si \emph{toda} la cadena coincide.
  \item \lstinline|regex_search(str, re)| — Comprueba si alguna \emph{subcadena} coincide.
  \item \lstinline|regex_replace(str, re, repl)| — Reemplaza coincidencias por otra cadena.
  \item \lstinline|sregex_iterator| — Permite recorrer todas las coincidencias.
\end{itemize}

\subsection{Ejemplos comunes}

\paragraph{1. Coincidencia completa}
\begin{lstlisting}[style=cppstyle]
regex re("[a-z]+");
cout << boolalpha << regex_match("abc", re) << endl; // true
cout << regex_match("abc123", re) << endl;           // false
\end{lstlisting}

\paragraph{2. Buscar dentro de una cadena}
\begin{lstlisting}[style=cppstyle]
string s = "code123forces";
regex re("[0-9]+");
smatch m;
if (regex_search(s, m, re))
    cout << "Primer numero: " << m.str() << "\n"; // "123"
\end{lstlisting}

\paragraph{3. Extraer todas las coincidencias}
\begin{lstlisting}[style=cppstyle]
string text = "ID12 AB34 XY99";
regex re("[A-Z]{2}[0-9]{2}");
for (sregex_iterator it(text.begin(), text.end(), re), end_it; it != end_it; ++it)
    cout << it->str() << " "; // ID12 AB34 XY99
\end{lstlisting}

\paragraph{4. Reemplazar coincidencias}
\begin{lstlisting}[style=cppstyle]
string s = "2025-10-17";
regex re("-"); 
string out = regex_replace(s, re, "/");
cout << out << "\n"; // 2025/10/17
\end{lstlisting}

\subsection{Sintaxis básica de patrones}

\begin{center}
\begin{tabular}{ll}
\hline
Patrón & Significado \\
\hline
\lstinline|.| & Cualquier carácter \\
\lstinline|[abc]| & Uno de a, b, c \\
\lstinline|[^abc]| & Cualquier excepto a, b, c \\
\lstinline|[a-z]| & Rango de letras minúsculas \\
\lstinline|\d| & Dígito (0–9) \\
\lstinline|\w| & Alfanumérico o \_ \\
\lstinline|\s| & Espacio, tab, salto de línea \\
\lstinline|^| & Inicio de cadena \\
\lstinline|$| & Fin de cadena \\
\lstinline|a?| & 0 o 1 ocurrencia de a \\
\lstinline|a*| & 0 o más ocurrencias \\
\lstinline|a+| & 1 o más ocurrencias \\
\lstinline|a{3}| & Exactamente 3 veces \\
\lstinline|a{2,4}| & Entre 2 y 4 veces \\
\lstinline|(abc)| & Grupo \\
\lstinline|a|b| & a o b \\
\hline
\end{tabular}
\end{center}

\subsection{Ejemplos útiles en práctica competitiva}

\paragraph{Validar número entero o decimal}
\begin{lstlisting}[style=cppstyle]
regex re("^-?[0-9]+(\\.[0-9]+)?$");
cout << regex_match("-12.34", re); // true
\end{lstlisting}

\paragraph{Detectar emails simples}
\begin{lstlisting}[style=cppstyle]
regex email("[\\w.-]+@[\\w.-]+\\.[a-zA-Z]{2,}");
cout << regex_match("name@example.com", email); // true
\end{lstlisting}

\paragraph{Separar palabras por espacios}
\begin{lstlisting}[style=cppstyle]
string s = "this is a test";
regex re("\\s+");
sregex_token_iterator it(s.begin(), s.end(), re, -1);
for (; it != sregex_token_iterator(); ++it)
    cout << *it << "\n";
\end{lstlisting}

\paragraph{Filtrar líneas que contengan un patrón}
\begin{lstlisting}[style=cppstyle]
regex pattern("ERROR");
string line;
while(getline(cin, line))
    if(regex_search(line, pattern))
        cout << line << "\n";
\end{lstlisting}

\subsection{Trucos y advertencias}

\begin{itemize}
  \item Para \textbf{competencias}, preferir \lstinline|string::find| o \lstinline|substr| si el patrón es simple: regex puede ser lento.
  \item Usar \lstinline|smatch| (para \lstinline|string|) o \lstinline|cmatch| (para C-strings).
  \item Los backslashes deben duplicarse: \lstinline|"\\d+"| (no \lstinline|"\d+"|).
  \item \lstinline|regex_constants::icase| hace coincidencias sin distinguir mayúsculas.
\end{itemize}

\subsection{Ejemplo completo: contar palabras con formato específico}
\begin{lstlisting}[style=cppstyle]
#include <bits/stdc++.h>
using namespace std;
int main(){
    string text = "A1 B22 C333 D44";
    regex re("[A-Z][0-9]{2,}");
    int cnt = 0;
    for(sregex_iterator it(text.begin(), text.end(), re), end_it; it!=end_it; ++it)
        cnt++;
    cout << "Coincidencias: " << cnt << "\n"; // 3
}
\end{lstlisting}

\subsection{Resumen de usos típicos}

\begin{center}
\begin{tabular}{ll}
\hline
Tarea & Regex típico \\
\hline
Número entero & \lstinline|^-?[0-9]+$| \\
Decimal & \lstinline|^-?[0-9]*\\.?[0-9]+$| \\
Email simple & \lstinline|[\\w.-]+@[\\w.-]+\\.[a-zA-Z]{2,}| \\
Solo letras & \lstinline|^[A-Za-z]+$| \\
Tokenizar por espacios & \lstinline|\\s+| \\
Buscar palabra específica & \lstinline|\\bword\\b| \\
\hline
\end{tabular}
\end{center}


\section{Matemáticas y Teoría de Números}

Conceptos y algoritmos esenciales para resolver problemas de aritmética, combinatoria y teoría de números en programación competitiva.

\subsection{Operaciones Básicas}

\begin{lstlisting}[style=cppstyle]
ll gcd(ll a, ll b){ return b ? gcd(b, a % b) : a; }
ll lcm(ll a, ll b){ return a / gcd(a,b) * b; }
ll mod(ll a, ll m){ return (a % m + m) % m; } // siempre positivo
\end{lstlisting}

\textbf{Trucos:}
\begin{itemize}
  \item Usar \lstinline|__gcd(a,b)| de \lstinline|<algorithm>|.
  \item Para overflow en multiplicación: usar \lstinline|(__int128)| o reducción modular.
\end{itemize}


\subsection{Aritmética Modular}

\paragraph{Suma, Resta, Multiplicación}
\begin{lstlisting}[style=cppstyle]
const ll MOD = 1e9+7;
ll add(ll a, ll b){ return (a + b) % MOD; }
ll sub(ll a, ll b){ return (a - b + MOD) % MOD; }
ll mul(ll a, ll b){ return (a * b) % MOD; }
\end{lstlisting}

\paragraph{Exponenciación Modular}
\begin{lstlisting}[style=cppstyle]
ll modpow(ll a, ll e, ll m){
    ll res = 1;
    while(e){
        if(e & 1) res = res * a % m;
        a = a * a % m;
        e >>= 1;
    }
    return res;
}
\end{lstlisting}

\paragraph{Inverso Modular (cuando \(m\) es primo)}
\begin{lstlisting}[style=cppstyle]
// Fermat: a^(m-2) mod m
ll modinv(ll a, ll m){ return modpow(a, m-2, m); }
\end{lstlisting}

\paragraph{Inversos del 1..n en O(n)}
\begin{lstlisting}[style=cppstyle]
vector<ll> inv(n+1);
inv[1] = 1;
for(int i=2; i<=n; i++)
    inv[i] = MOD - (MOD/i) * inv[MOD % i] % MOD;
\end{lstlisting}


\subsection{Números Primos}

\paragraph{Prueba de primalidad simple}
\begin{lstlisting}[style=cppstyle]
bool is_prime(ll n){
    if(n < 2) return false;
    for(ll i=2; i*i<=n; i++)
        if(n % i == 0) return false;
    return true;
}
\end{lstlisting}

\paragraph{Criba de Eratóstenes}
\begin{lstlisting}[style=cppstyle]
int N = 1e6;
vector<int> primes;
vector<bool> isprime(N+1, true);

void sieve(){
    isprime[0] = isprime[1] = false;
    for(int i=2; i*i<=N; i++)
        if(isprime[i])
            for(int j=i*i; j<=N; j+=i)
                isprime[j] = false;
    for(int i=2; i<=N; i++)
        if(isprime[i]) primes.push_back(i);
}
\end{lstlisting}

\paragraph{Factorización con Criba}
\begin{lstlisting}[style=cppstyle]
vector<int> spf(N+1); // smallest prime factor

void sieve_spf(){
    iota(spf.begin(), spf.end(), 0);
    for(int i=2; i*i<=N; i++)
        if(spf[i] == i)
            for(int j=i*i; j<=N; j+=i)
                if(spf[j] == j) spf[j] = i;
}

vector<int> factorize(int n){
    vector<int> f;
    while(n > 1){
        f.push_back(spf[n]);
        n /= spf[n];
    }
    return f;
}
\end{lstlisting}


\subsection{Combinatoria y Factoriales Modulares}

\paragraph{Precomputar factoriales}
\begin{lstlisting}[style=cppstyle]
const int MAXN = 1e6;
vector<ll> fact(MAXN+1), invfact(MAXN+1);

void init_fact(){
    fact[0] = 1;
    for(int i=1;i<=MAXN;i++) fact[i] = fact[i-1]*i % MOD;
    invfact[MAXN] = modinv(fact[MAXN], MOD);
    for(int i=MAXN;i>=1;i--) invfact[i-1] = invfact[i]*i % MOD;
}
\end{lstlisting}

\paragraph{Combinaciones nCr mod m}
\begin{lstlisting}[style=cppstyle]
ll nCr(int n, int r){
    if(r<0 || r>n) return 0;
    return fact[n] * invfact[r] % MOD * invfact[n-r] % MOD;
}
\end{lstlisting}

\textbf{Truco:}  
Para \(n \le 10^6\), precomputar factoriales; para \(n\) grande pero \(r\) pequeño, usar fórmula multiplicativa.


\subsection{Algoritmos Numéricos Clásicos}

\paragraph{MCD extendido}
\begin{lstlisting}[style=cppstyle]
ll extgcd(ll a, ll b, ll &x, ll &y){
    if(b == 0){ x = 1; y = 0; return a; }
    ll x1, y1;
    ll g = extgcd(b, a % b, x1, y1);
    x = y1;
    y = x1 - y1 * (a / b);
    return g;
}
\end{lstlisting}

\paragraph{Solución de Congruencia Lineal \(a x \equiv b \pmod{m}\)}
\begin{lstlisting}[style=cppstyle]
ll solve_congruence(ll a, ll b, ll m){
    ll x, y;
    ll g = extgcd(a, m, x, y);
    if(b % g) return -1; // no solution
    x = (x * (b/g)) % m;
    return (x + m) % m;
}
\end{lstlisting}

\paragraph{Criba para funciones aritméticas (\(\phi\))}
\begin{lstlisting}[style=cppstyle]
vector<int> phi(N+1);
void sieve_phi(){
    for(int i=0;i<=N;i++) phi[i]=i;
    for(int i=2;i<=N;i++)
        if(phi[i]==i)
            for(int j=i;j<=N;j+=i)
                phi[j]-=phi[j]/i;
}
\end{lstlisting}

\paragraph{Conteo de divisores y suma de divisores}
\begin{lstlisting}[style=cppstyle]
int num_div(int n){
    int cnt=1;
    for(int i=2, c;i*i<=n;i++){
        c=0;
        while(n%i==0) n/=i, c++;
        cnt *= (c+1);
    }
    if(n>1) cnt*=2;
    return cnt;
}
\end{lstlisting}


\subsection{Trucos y Consejos}

\begin{itemize}
  \item \textbf{MODs comunes:} \(10^9+7\), \(998244353\).
  \item Siempre usar \lstinline|(a * 1LL * b) % MOD| para evitar overflow.
  \item Para \textbf{coprimos grandes}, el inverso modular existe.
  \item Para \textbf{conteo combinatorio}, usar logaritmos o Lucas si \(n\) es muy grande.
  \item \textbf{Euler’s Totient:} \(\phi(n)\) = cantidad de números \(\le n\) coprimos con \(n\).
  \item \textbf{Exponenciación rápida} = base para muchos algoritmos (potencias, inversos, etc.).
\end{itemize}


\subsection{Resumen de Complejidades}

\begin{center}
\begin{tabular}{lcc}
\hline
Algoritmo & Complejidad & Descripción \\
\hline
Criba Eratóstenes & \(O(n \log\log n)\) & Todos los primos hasta \(n\) \\
Factorización (con SPF) & \(O(\log n)\) & Usando preprocesamiento \\
Exp. Modular & \(O(\log e)\) & Potencia rápida \\
Inversos 1..n & \(O(n)\) & Todos los inversos mod primo \\
nCr con factoriales & \(O(1)\) & Tras precomputar factoriales \\
\hline
\end{tabular}
\end{center}




\end{document}
